\documentclass{article}

\usepackage{iclr2026_conference,times}

\input{math_commands.tex}

\usepackage{amsmath}
\usepackage{amssymb}
\usepackage{array}
\usepackage{booktabs}
\usepackage{xcolor}
\usepackage{hyperref}
\usepackage{url}

\newcommand{\method}{HSE}
\newcommand{\skillrepo}{SkillRepo}
\newcommand{\todo}[1]{\textcolor{red}{[TODO: #1]}}
\newcommand{\tbd}{\textcolor{red}{TBD}}
\newcolumntype{L}[1]{>{\raggedright\arraybackslash}p{#1}}
\newcolumntype{C}[1]{>{\centering\arraybackslash}p{#1}}

\title{Verifier-Grounded Hierarchical Skill Evolution\\for Software Engineering Agents}

% Anonymous submission. Uncomment \iclrfinalcopy only for camera-ready.
\author{Anonymous authors}

% \iclrfinalcopy
\begin{document}

\maketitle

\begin{abstract}
Software-engineering agents leave behind rich traces: issue interpretations, file
searches, edits, tests, failures, and recoveries. Turning these traces into
reusable knowledge is attractive, but raw trace retrieval risks task leakage and
one-shot summaries can introduce harmful or over-specific instructions. We
propose \method{}, a verifier-grounded framework that evolves external textual
skills for frozen coding agents. \method{} separates single-task evidence from
reusable skills, promotes evidence through a hierarchy of repo-level and
cross-repo failure-mode skills, and trains separate Skill Writer and Skill
Evaluator policies under a fixed harness policy. During SWEGym training, the
official verifier is used only for outcome supervision, evaluator calibration,
and validation gates. During downstream evaluation, verifier labels are not
available for updating memory. We describe a strict no-exact-transfer protocol
and an online test-time evolution protocol over a $2{\times}2$ executor matrix:
GLM-5.2 and GPT-5.5, each run in Pi and Claude-style agent harnesses. This
matrix lets us test model transfer, harness transfer, and joint model+harness
transfer. Final resolved rates and ablations are intentionally left blank in
this draft pending the full experimental run.
\end{abstract}

\section{Introduction}

Large language model agents can now repair real software bugs by reading a
repository, editing code, and running tests. Yet each successful or failed
attempt produces information that is usually discarded after the benchmark
trial. A natural next step is to convert agent trajectories into external
skills: compact textual procedures that tell future agents how to localize,
edit, validate, or recover on similar software tasks.

This conversion is not straightforward. A raw trajectory can contain
task-specific paths, accidental patch shapes, transient environment errors, and
implicit information about the benchmark instance. If such details are injected
into a future prompt, the system may improve by memorization rather than by
transfer. Conversely, a generic summary such as ``run focused tests'' is often
too weak to change behavior. The central question of this paper is therefore:
\emph{how can a coding-agent system learn reusable software-engineering skills
from training-time verifier feedback without using test-time verifier feedback
or leaking task-specific solutions?}

We propose \method{}: \textbf{hierarchical skill evolution}. The key design
choice is to not call every trace artifact a skill. A task attempt first becomes
\emph{task evidence}: a structured record of public signals such as touched
paths, edited files, tests, diagnostics, selected skills, failure signatures,
and verifier outcomes. A Skill Writer policy then drafts candidate repo-level
skills only from multiple task evidence records in the same repository. A
candidate failure-mode skill is drafted only from repeated patterns across
repositories, such as localization drift, weak validation, no-diff exits, or
timeout recovery. A Skill Evaluator policy scores each candidate for utility,
support, and risk, and decides whether to accept, reject, revise, or keep it as
memory-only evidence.

The framework is built around a hard information boundary. SWEGym verifiers are
available during training, so we use them to label outcomes, calibrate the
evaluator, and gate candidate skill states on held-out validation tasks. On
SWE-Bench Verified and other downstream benchmarks, no official verifier label,
hidden test result, gold patch, or future task outcome may influence skill
updates. This makes the resulting protocol suitable for testing whether
SWEGym-trained skills and writer/evaluator policies transfer to new software
tasks.

Our contributions are:
\begin{itemize}
    \item We formulate software-agent skill evolution as a hierarchical
    evidence-gated process that separates task evidence, repo-level skills,
    cross-repo failure-mode skills, and writer/evaluator policies.
    \item We introduce a verifier-boundary protocol: training may use SWEGym
    verifier outcomes, while downstream test-time evolution is gated only by a
    frozen evaluator over public trace signals.
    \item We give an implementation for the Pi coding-agent harness with an
    OpenAI-compatible model interface, Harbor/E2B execution, versioned skill
    archives, promotion logs, validation gates, and deterministic evaluator
    shells that can be replaced by learned policies.
    \item We specify the ICLR experimental protocol for GLM-5.2 and GPT-5.5 in
    both Pi and Claude-style harnesses, including model-transfer and
    harness-transfer tests, with exact-task memory reported only as an
    upper-bound adaptation setting.
\end{itemize}

\section{Problem Setting}

\paragraph{Tasks and agents.}
Let $x$ be a software-engineering task consisting of a natural-language issue,
a repository snapshot, and a public execution environment. A frozen coding
agent $M$ interacts with a harness $H$ by issuing file reads, searches, edits,
shell commands, and final answers. A trajectory $\tau$ contains the observable
history: prompts, tool calls, outputs, diffs, public tests, and runtime
diagnostics.

\paragraph{Training-time verifier.}
During SWEGym training we have a verifier $V$ that maps a completed trial to a
reward $r \in \{0,1\}$ or a diagnostic outcome. We use $V$ as supervision for
training-time labels and validation gates, but not as a writer of skills and not
as a provider of hidden solutions. At downstream test time, $V$ is unavailable
for memory updates. It is used only after the protocol completes to report final
benchmark scores.

\paragraph{Skill state.}
The external state $S$ contains structured evidence and accepted skills. We use
four deployed categories:
\begin{itemize}
    \item \textbf{Task evidence} records what happened on one task. It is not a
    reusable skill.
    \item \textbf{Repo skills} summarize repeated, evidence-gated procedures
    within a repository.
    \item \textbf{Failure-mode skills} summarize cross-repo recovery strategies
    for repeated agent failures.
    \item \textbf{General seed skills} provide weak process priors for common
    stages such as reproduce, localize, edit, validate, and recover.
\end{itemize}
The writer policy $\pi_w$ drafts evidence summaries and candidate skills. The
evaluator policy $\pi_e$ assigns decisions, proxy rewards, confidence, and risk
flags.

\section{Hierarchical Skill Evolution}

\subsection{State Hierarchy}

\method{} uses different update clocks for different abstraction levels. Lower
levels can change quickly because they affect only local evidence. Higher levels
change slowly because they can influence future tasks.

\begin{table}[t]
\caption{Skill-evolution objects and update clocks. Main experiments freeze all
training-derived transferable state before downstream evaluation.}
\label{tab:levels}
\begin{center}
\begin{tabular}{L{0.12\linewidth} L{0.25\linewidth} L{0.23\linewidth} L{0.25\linewidth}}
\toprule
Level & Object & Update clock & Evidence requirement \\
\midrule
L0 & Rollout trace & Every task attempt & Public trajectory and runtime outputs \\
L1 & Task evidence & Every task attempt & One task trace; verifier label during training only \\
L2 & Repo skill candidate & Same-repo batch & Multiple task evidence records from one repo \\
L3 & Failure-mode skill candidate & Cross-repo trigger & Repeated repo-level summaries across repos \\
L4 & Skill Writer policy & Slow policy window & Aggregated accepted/rejected histories \\
L5 & Skill Evaluator policy & Slow policy window & Verifier-calibrated false accepts/rejects \\
\bottomrule
\end{tabular}
\end{center}
\end{table}

Task evidence stores public fields such as task id, repo, problem summary,
reward, previous reward, case label, exception type, selected skills, touched
paths, edited paths, public test commands, failure signature, and diagnostic
signals. It may feed higher-level aggregation, but it is not rendered as a
reusable \texttt{SKILL.md} artifact by default.

\subsection{Writer and Evaluator}

The writer policy produces structured candidates. For a repo candidate, the
writer receives a cluster of same-repo task evidence and emits the fields
name, trigger, evidence gate, actions, validation hint, abort condition, and
support summary. For a failure-mode candidate, it receives a cross-repo cluster
and emits a process-level recovery skill with a failure signature, trigger,
actions, stop condition, and support summary.
\[
c_{\mathrm{repo}} = f_w(E_{\mathrm{repo}}), \qquad
c_{\mathrm{fail}} = f_w(E_{\mathrm{cross\text{-}repo}}).
\]
Both candidate types are deliberately structured so the evaluator can inspect
their triggers, evidence support, actions, and risks before deployment.

The evaluator implements the boundary between candidate generation and deployed
skills:
\[
\pi_e(c, E) \rightarrow
(d, \hat r, \rho, z),
\]
where $d \in \{\textsc{accept}, \textsc{reject}, \textsc{revise},
\textsc{memory-only}\}$, $\hat r$ is a proxy reward, $\rho$ is confidence, and
$z$ is a set of risk flags. Hard safety risks include hidden verifier details,
oracle patches, exact patches, task-id memorization, and private sandbox paths.
Repo candidates require repeated same-repo support; failure-mode candidates
require support from multiple distinct repositories.

Our current implementation uses a deterministic evaluator shell with explicit
thresholds and risk checks. This is a deliberate interface choice: it makes the
contract auditable now, while allowing the same function call to be replaced by
a learned evaluator policy trained on SWEGym calibration events.

\begin{figure}[t]
\begin{center}
\fbox{\begin{minipage}{0.93\linewidth}
\small
\textbf{Agent trace}
$\rightarrow$ \textbf{Task evidence}
$\rightarrow$ \textbf{Repo candidate}
$\rightarrow$ \textbf{Failure-mode skill}.
The trace contains tool calls, edits, tests, and diagnostics. Task evidence is
a public structured record, not a deployed transferable skill. Repo candidates
summarize same-repo repeated paths, tests, and failure signals. Failure-mode
skills summarize cross-repo recovery rules with triggers and stop conditions.
\vspace{0.5em}
\textbf{Evaluator gate:} each promotion step must pass support thresholds,
risk checks, and validation constraints before it becomes active in the
\skillrepo{}.
\end{minipage}}
\end{center}
\caption{\textbf{AI-friendly caption:} A left-to-right pipeline showing how a
software-engineering agent trace is converted into task evidence, then into a
repo-level candidate, then into a cross-repo failure-mode skill. The visual
should emphasize that task evidence is not directly deployed as a transferable
skill; every upward arrow is guarded by an evaluator gate that checks support,
risk, and validation.}
\label{fig:pipeline}
\end{figure}

\subsection{Fixed Harness Policy}

The learned writer and evaluator are not allowed to invent their own reward
labels, update clocks, or promotion budgets. A fixed harness policy maps
verifier transitions and runtime diagnostics to case labels:
\begin{itemize}
    \item \textbf{Strong positive}: previous attempt failed and current attempt
    passed.
    \item \textbf{Weak positive}: current attempt passed without a paired
    improvement, or a previously passing attempt remains passing.
    \item \textbf{Strong negative}: previous attempt passed and current attempt
    failed.
    \item \textbf{Weak negative}: current attempt failed without a paired
    regression.
    \item \textbf{Diagnostic}: missing reward, timeout, setup error, verifier
    error, sandbox failure, or provider/runtime failure.
\end{itemize}

Cold-start solved traces are weak positives; unsolved traces are weak
negatives unless a diagnostic is present. Strong positives can increase support
for a candidate, but cannot directly create a global skill. Strong negatives
prioritize rollback, trigger narrowing, and evaluator false-accept penalties.
Diagnostics are kept out of semantic edit skills and are used only for runtime
or recover/validate safeguards.

\subsection{Training Loop}

The training loop alternates between rollouts, evidence extraction, candidate
generation, evaluator gating, and validation.

\begin{figure}[t]
\begin{center}
\fbox{\begin{minipage}{0.94\linewidth}
\small
\textbf{Training on SWEGym:}
\[
\text{rollout} \rightarrow \text{task evidence} \rightarrow
\text{repo/failure candidates} \rightarrow
\pi_e\text{ gate} \rightarrow
\text{validation gate}.
\]
The SWEGym verifier labels outcomes and validates candidate skill states.

\vspace{0.5em}
\textbf{Downstream evaluation:}
\[
\text{frozen }(\pi_w,\pi_e,\text{general/failure skills})
\rightarrow \text{agent rollout} \rightarrow
\text{public evidence only}.
\]
No official downstream verifier label, hidden test, gold patch, or future task
outcome is used for updating skills.
\end{minipage}}
\end{center}
\caption{\textbf{AI-friendly caption:} A two-lane diagram contrasting training
and downstream evaluation. The top lane is SWEGym training, where verifier
labels supervise case labels, evaluator calibration, and validation gates. The
bottom lane is downstream testing, where the writer, evaluator, and transferable
skills are frozen and only public trace signals can be used. The verifier icon
must appear only in the training lane and final reporting, not in the update
path for downstream tests.}
\label{fig:verifier-boundary}
\end{figure}

\noindent\textbf{Default update hyperparameters.}
Unless otherwise stated, repo updates use same-repo batches of 5 task evidence
events, failure-mode updates trigger every 4 repo updates, failure-mode
promotion requires at least 3 supporting repositories, and policy updates are
limited to at most 3 bullet edits per policy window. Each repo batch emits at
most one repo candidate by default; each failure-mode trigger promotes at most
two accepted failure-mode skills. These values match the current harness
defaults and are treated as fixed protocol hyperparameters.

\subsection{Validation Gate}

Candidate skill states are evaluated on a repo-isolated SWEGym validation set.
We compute:
\[
\mathrm{effective\_resolved\_rate}
=
\frac{\mathrm{resolved\ valid\ trials}}{\mathrm{effective\ valid\ trials}},
\quad
\mathrm{diagnostic\_rate}
=
\frac{\mathrm{diagnostic\ trials}}{\mathrm{total\ trials}}.
\]
A candidate state is accepted only if it has enough effective validation trials,
does not exceed the diagnostic-rate threshold, and does not regress relative to
the current accepted state beyond a fixed tolerance. Otherwise, the candidate
version is rejected or rolled back, and the decision is logged.

\section{Experiments}

This section specifies the experiment settings. Numeric results are reserved
until the full ICLR run finishes.

\subsection{Datasets}

\paragraph{SWEGym training and validation.}
We use a 500-task SWEGym subset as the evolution source. The split is
repo-isolated: repositories selected for validation are held out from the
training split, with the validation size chosen closest to 5\% of the subset.
In the planned main run this yields 474 training tasks and 26 validation tasks.
The SWEGym verifier is available only in this phase.

\paragraph{SWE-Bench Verified final evaluation.}
The primary downstream benchmark is SWE-Bench Verified. No exact SWE-Bench task
trajectory, official verifier label, hidden test output, or gold patch is used
to update skills during the evaluated run. The official verifier is invoked
only after the run to report resolved count and resolved rate.

\paragraph{Optional out-of-distribution evaluation.}
We will also report transfer on SWE-Bench Pro and Terminal-Bench 2.0 when the
runs are complete. Terminal-Bench uses the same Harbor/E2B execution wrapper
but a terminal-task prompt.

\subsection{Models and Harnesses}

\paragraph{Executor matrix.}
We evaluate two executor models, GLM-5.2 (\texttt{glm52}) and GPT-5.5, each
under two agent harnesses: Pi Coding Agent and a Claude-style coding-agent
harness. All four cells use the same Harbor/E2B sandbox interface, benchmark
task selection, task ordering, timeout budget, validation gate, and
no-test-verifier update rule. Where model APIs differ, we use the same
OpenAI-compatible provider abstraction and record the endpoint type in the run
config.

\paragraph{Model and harness transfer.}
The $2{\times}2$ matrix lets us distinguish three transfer questions. First,
\textbf{model transfer}: skills and writer/evaluator policies trained with one
model are evaluated with the other model under the same harness
(GLM-5.2$\rightarrow$GPT-5.5 and GPT-5.5$\rightarrow$GLM-5.2). Second,
\textbf{harness transfer}: skills trained under one harness are evaluated under
the other harness with the same model (Pi$\rightarrow$Claude-style and
Claude-style$\rightarrow$Pi). Third, \textbf{joint transfer}: both the model and
harness change between evolution and evaluation. We report these separately
from within-cell evolution, where train and test use the same model/harness
cell.

\paragraph{Tools and controls.}
The agent tool surface is fixed within each harness family and exposes the same
software-engineering operations: read, write, edit, bash, grep, find, and ls.
Prompts may differ between Pi and Claude-style harnesses, but task order,
sandbox resources, retry policy, and acceptance rules are held fixed. Each run
logs the active skill version, source model/harness cell, target model/harness
cell, selected skills, runtime diagnostics, verifier reports, and score
reports.

\paragraph{Execution controls.}
All compared settings use the same dataset version, task order, sandbox
resources, agent timeout, retry policy, and concurrency within a model/harness
pair. We log job directories, configs, active skill versions, selected skills,
tool traces, runtime diagnostics, verifier reports, and score reports.

\subsection{Protocols}

\begin{table}[t]
\caption{Main evaluation protocols. Exact-task skills are reported only as an
adaptation upper bound and are not used to support transfer claims.}
\label{tab:protocols}
\begin{center}
\begin{tabular}{L{0.25\linewidth} C{0.15\linewidth} C{0.16\linewidth} L{0.32\linewidth}}
\toprule
Setting & Uses exact task memory? & Updates during test? & Purpose \\
\midrule
No Skills & No & No & Frozen-agent baseline \\
Public SWE Skills & No & No & Static public skill baseline \\
Raw Trace Retrieval & No & No & Tests whether raw trajectory memory is enough \\
One-shot StageSkill & No & No & Tests ungated trace-to-skill summarization \\
Frozen Failure Skills & No & No & Tests SWEGym-trained transferable failure-mode skills \\
Frozen Writer/Evaluator + Online Repo Skills & No & Yes, public evidence only & Tests verifier-free online evolution \\
Same-budget Retry & No & No structured skills & Controls for extra attempts or compute \\
Exact Task Skills & Yes & Optional & Upper bound/adaptation only \\
\bottomrule
\end{tabular}
\end{center}
\end{table}

\subsection{Metrics}

The primary metric is resolved rate. We also report resolved-count deltas,
paired transition counts ($0{\rightarrow}1$, $1{\rightarrow}0$,
$1{\rightarrow}1$, $0{\rightarrow}0$), diagnostic rate, no-diff rate, timeout
rate, weak-validation rate, localization-drift rate, selected-skill read rate,
trigger-match rate, retrieval false-positive rate, tool-call count, token count,
wall-clock time, and validation-command count. These metrics are required
because a skill system can improve pass rate while increasing regressions,
latency, or prompt cost.

\subsection{Result Tables to Fill}

\begin{table}[h]
\caption{Primary SWE-Bench Verified within-cell results. Values are
intentionally blank until the full GLM-5.2/GPT-5.5 by Pi/Claude-harness matrix
finishes.}
\label{tab:main-results}
\begin{center}
\begin{tabular}{L{0.27\linewidth} C{0.10\linewidth} C{0.10\linewidth} C{0.11\linewidth} C{0.11\linewidth} C{0.11\linewidth}}
\toprule
Model/Harness/Setting & Resolved & Rate & $0{\rightarrow}1$ & $1{\rightarrow}0$ & Diag. \\
\midrule
GLM-5.2 / Pi / No Skills & \tbd & \tbd & \tbd & \tbd & \tbd \\
GLM-5.2 / Pi / \method{} & \tbd & \tbd & \tbd & \tbd & \tbd \\
GLM-5.2 / Claude harness / No Skills & \tbd & \tbd & \tbd & \tbd & \tbd \\
GLM-5.2 / Claude harness / \method{} & \tbd & \tbd & \tbd & \tbd & \tbd \\
GPT-5.5 / Pi / No Skills & \tbd & \tbd & \tbd & \tbd & \tbd \\
GPT-5.5 / Pi / \method{} & \tbd & \tbd & \tbd & \tbd & \tbd \\
GPT-5.5 / Claude harness / No Skills & \tbd & \tbd & \tbd & \tbd & \tbd \\
GPT-5.5 / Claude harness / \method{} & \tbd & \tbd & \tbd & \tbd & \tbd \\
\bottomrule
\end{tabular}
\end{center}
\end{table}

\begin{table}[h]
\caption{Transfer matrix to fill after the full run. Source denotes the
model/harness cell used for skill evolution; target denotes the cell used for
SWE-Bench Verified evaluation.}
\label{tab:transfer}
\begin{center}
\begin{tabular}{L{0.25\linewidth} L{0.25\linewidth} L{0.20\linewidth} C{0.13\linewidth}}
\toprule
Source cell & Target cell & Transfer type & Resolved delta \\
\midrule
GLM-5.2 / Pi & GPT-5.5 / Pi & Model transfer & \tbd \\
GPT-5.5 / Pi & GLM-5.2 / Pi & Model transfer & \tbd \\
GLM-5.2 / Claude harness & GPT-5.5 / Claude harness & Model transfer & \tbd \\
GPT-5.5 / Claude harness & GLM-5.2 / Claude harness & Model transfer & \tbd \\
GLM-5.2 / Pi & GLM-5.2 / Claude harness & Harness transfer & \tbd \\
GLM-5.2 / Claude harness & GLM-5.2 / Pi & Harness transfer & \tbd \\
GPT-5.5 / Pi & GPT-5.5 / Claude harness & Harness transfer & \tbd \\
GPT-5.5 / Claude harness & GPT-5.5 / Pi & Harness transfer & \tbd \\
GLM-5.2 / Pi & GPT-5.5 / Claude harness & Joint transfer & \tbd \\
GPT-5.5 / Claude harness & GLM-5.2 / Pi & Joint transfer & \tbd \\
\bottomrule
\end{tabular}
\end{center}
\end{table}

\begin{table}[h]
\caption{Ablation plan. Each row removes one mechanism while keeping model,
harness, task order, and budget fixed.}
\label{tab:ablations}
\begin{center}
\begin{tabular}{L{0.30\linewidth} L{0.42\linewidth} C{0.18\linewidth}}
\toprule
Ablation & Question & Resolved delta \\
\midrule
Without evaluator gate & Do ungated candidates introduce harmful skills? & \tbd \\
Without validation gate & Does repo-isolated validation prevent regressions? & \tbd \\
Without failure-mode skills & Are cross-repo recovery rules useful? & \tbd \\
Without repo skills & Are same-repo procedures useful in online evolution? & \tbd \\
Without policy update windows & Do slow clocks reduce drift and prompt bloat? & \tbd \\
Raw trace retrieval only & Is structured skill abstraction necessary? & \tbd \\
\bottomrule
\end{tabular}
\end{center}
\end{table}

\section{Related Work}

\paragraph{Software-engineering benchmarks.}
SWE-Bench introduced real GitHub issue resolution as a benchmark for language
models and coding agents \citep{swebench2024}. SWE-Bench Verified refines the
benchmark with human-verified tasks and evaluation criteria
\citep{swebenchverified2024}. SWEGym provides executable training environments
for software-engineering agents and is the source of verifier-supervised skill
evolution in our work \citep{swegym2025}. Terminal-Bench studies broader
terminal workflows beyond repository patching \citep{terminalbench2025}.

\paragraph{Language-agent memory and reflection.}
Prior language-agent systems use verbal feedback, self-reflection, or memory to
improve future behavior \citep{react2023,reflexion2023,voyager2023}. \method{}
differs by enforcing a benchmark-verifier boundary and by separating
task-local evidence from transferable repo and failure-mode skills.

\paragraph{Skill extraction and text-space optimization.}
Recent systems distill trajectories into textual skills, optimize external
instructions, or maintain skill repositories for agents
\citep{trace2skill2024,skillopt2024,skillos2025,autoskill2023}. Our focus is
the software-engineering setting, where executable verifiers are available
during training but not during test-time memory updates, and where exact task
leakage is a central evaluation risk.

\section{Limitations}

This draft intentionally leaves experimental results blank. The method also has
several substantive limitations. First, verifier-grounded training is only as
good as the verifier and validation split; incomplete public tests can still
miscalibrate evaluator decisions. Second, failure-mode skills may be too broad
if failure signatures are noisy. Third, online test-time evolution is order
sensitive, so we must report original-order and random-order robustness. Fourth,
the current evaluator implementation is deterministic and rule-based; a learned
evaluator must be tested against false accepts, false rejects, and prompt bloat.
Finally, the full $2{\times}2$ transfer matrix is a planned setting. Within-cell
results, model-transfer results, harness-transfer results, and joint-transfer
results must be reported separately because they answer different questions.

\section{Conclusion}

We presented \method{}, a verifier-grounded approach for evolving external
software-engineering skills for frozen coding agents. The framework promotes
public task evidence through repo-level and cross-repo failure-mode gates,
separates writer and evaluator policies, and freezes transferable state before
downstream evaluation. The resulting protocol is designed to answer whether
SWEGym can train reusable skill-writing and skill-selection behavior that
transfers to SWE-Bench Verified, SWE-Bench Pro, and Terminal-Bench without
test-time verifier feedback.

\bibliographystyle{iclr2026_conference}
\bibliography{references}

\appendix
\section{Implementation Notes}

The implementation stores accepted skills under versioned directories
\texttt{skills/accepted/vXXXX/}. Repo skills are written under
\texttt{\_repos/}, failure-mode skills under \texttt{\_failure\_modes/}, and
general SWE seeds under \texttt{\_general/swe/}. Training artifacts include
\texttt{task\_events.jsonl}, \texttt{repo\_clusters.jsonl},
\texttt{failure\_clusters.jsonl}, \texttt{evaluator\_decisions.jsonl},
\texttt{evaluator\_calibration.jsonl}, \texttt{promotion\_decisions.jsonl},
\texttt{harness\_policy.md}, \texttt{generator\_policy.md}, and
\texttt{evaluator\_policy.md}.

\section{Command Skeletons}

\begin{verbatim}
# SWEGym transfer loop, one source cell.
python scripts/run_swegym_skill_evo_loop.py \
  --run-name swegym_<model>_<harness>_hse_main \
  --provider openai \
  --provider-model <model-id> \
  --provider-api openai-completions \
  --concurrency <N>

# SWE-Bench Verified target-cell evaluation with accepted skills.
python scripts/run_skill_evo_verified.py \
  --run-name verified_<source>_to_<target>_hse_final \
  --dataset swe-bench/swe-bench-verified@2 \
  --provider openai \
  --provider-model <target-model-id> \
  --all-verified \
  --concurrency <N>
\end{verbatim}

\end{document}
